\documentclass[../main.tex]{subfiles}

\begin{document}

\chapter*{Capaian Pembelajaran Buku}
\addcontentsline{toc}{chapter}{Capaian Pembelajaran Buku}

Bagian ini merangkum capaian pembelajaran yang dibebankan pada buku ajar ini
beserta pemetaannya ke bab-bab. Rumusan CPL, CPMK, dan Sub-CPMK mengikuti
silabus mata kuliah Kriptografi yang berbasis \textit{Outcome-Based Education}
(OBE).

\section*{Capaian Pembelajaran Lulusan (CPL)}

\begin{itemize}
    \item \textbf{PP1 (Penguasaan Pengetahuan):} Menguasai konsep teoritis dan
          prinsip-prinsip dasar dalam matematika murni dan terapan untuk
          menyelesaikan berbagai permasalahan.
    \item \textbf{KK1 (Keterampilan Khusus):} Mampu menganalisis pemikiran
          matematis, pemecahan masalah, serta model matematika dengan pendekatan
          matematis dan teknologi untuk menghasilkan interpretasi yang akurat.
    \item \textbf{KU2 (Keterampilan Umum):} Mampu mengkaji dan menerapkan ilmu
          pengetahuan serta teknologi dengan menyusun deskripsi saintifik serta
          dapat dipertanggungjawabkan.
    \item \textbf{KU3 (Keterampilan Umum):} Mampu mendokumentasikan, menyimpan,
          mengamankan, dan menemukan kembali data untuk menjamin kesahihan dan
          mencegah plagiasi.
    \item \textbf{S2 (Sikap):} Bertindak dengan tanggung jawab, kepedulian, dan
          etika dalam kehidupan bermasyarakat, berbangsa, dan bernegara
          berdasarkan nilai kemanusiaan, kebhinekaan, dan Pancasila.
\end{itemize}

\section*{Capaian Pembelajaran Mata Kuliah (CPMK)}

Setelah menyelesaikan mata kuliah ini, mahasiswa diharapkan mampu:

\begin{itemize}
    \item \textbf{CPMK-1} (dibebankan pada CPL: PP1, KK1, KU2): Mengenal
          beberapa jenis algoritma kriptografi klasik dan modern.
    \item \textbf{CPMK-2} (dibebankan pada CPL: PP1, KU2, KU3, S2): Memahami
          konsep kriptografi secara umum dan urgensinya dalam dunia teknologi
          informasi.
\end{itemize}

\section*{Sub-CPMK}

\noindent\textbf{CPMK-2 (konsep dan urgensi):}
\begin{itemize}
    \item \textbf{Sub-CPMK 2.1:} Menjelaskan definisi, terminologi, sejarah, dan
          layanan keamanan (kerahasiaan, integritas, ketersediaan,
          nir-penyangkalan).
    \item \textbf{Sub-CPMK 2.2:} Menjelaskan urgensi kriptografi dalam teknologi
          informasi (privasi data, transaksi elektronik, HTTPS/TLS, etika
          penggunaan).
    \item \textbf{Sub-CPMK 2.3:} Menjelaskan landasan matematika dasar yang
          relevan (aritmetika modulo, bilangan prima, invers modulo) secara
          pengenalan.
\end{itemize}

\noindent\textbf{CPMK-1 (algoritma klasik dan modern):}
\begin{itemize}
    \item \textbf{Sub-CPMK 1.1:} Mengenal cipher klasik (substitusi, transposisi,
          Vigen\`ere, Playfair, Affine, Hill, OTP) dan gagasan kriptanalisis
          frekuensi.
    \item \textbf{Sub-CPMK 1.2:} Mengenal kriptografi modern simetris
          (stream/block, Feistel, DES, AES, mode operasi dasar).
    \item \textbf{Sub-CPMK 1.3:} Mengenal kriptografi modern asimetris (RSA,
          Diffie--Hellman, ElGamal, pengantar ECC).
    \item \textbf{Sub-CPMK 1.4:} Mengenal fungsi hash, MAC, tanda tangan digital,
          dan peran protokol (SSL/TLS, PKI) sebagai aplikasi algoritma modern.
\end{itemize}

\section*{Pemetaan Sub-CPMK ke Bab}

Tabel berikut menunjukkan bab mana yang menjadi wahana pencapaian tiap
Sub-CPMK.

\begin{longtable}{|c|>{\raggedright\arraybackslash}p{4.2cm}|>{\raggedright\arraybackslash}p{6.0cm}|}
\hline
\textbf{Sub-CPMK} & \textbf{CPMK} & \textbf{Bab} \\
\hline
\endfirsthead
\hline
\textbf{Sub-CPMK} & \textbf{CPMK} & \textbf{Bab} \\
\hline
\endhead
2.1 & CPMK-2 & Bab 1, Bab 2 \\
\hline
2.2 & CPMK-2 & Bab 1, Bab 2, Bab 16 \\
\hline
2.3 & CPMK-2 & Bab 3, Bab 13 \\
\hline
1.1 & CPMK-1 & Bab 4 \\
\hline
1.2 & CPMK-1 & Bab 5, Bab 6, Bab 7, Bab 8, Bab 9 \\
\hline
1.3 & CPMK-1 & Bab 10, Bab 11, Bab 12, Bab 13, Bab 14 \\
\hline
1.4 & CPMK-1 & Bab 15, Bab 16 \\
\hline
\end{longtable}

\textit{Catatan:} Bab 9 dan Bab 13 bersifat pengayaan, yaitu memperluas cakupan
CPMK-1 di luar daftar materi inti silabus. Rincian pemetaan minggu pertemuan ke
bab disajikan pada bagian \emph{Pemetaan Silabus terhadap Bab Buku}.

\end{document}
